\documentclass[11pt,a4paper]{article}
\usepackage[utf8]{inputenc}
\usepackage[T1]{fontenc}
\usepackage{tgpagella}
\usepackage[spanish]{babel}
\usepackage{geometry}
\usepackage{booktabs}
\usepackage{xcolor}
\usepackage{titlesec}
\usepackage{fancyhdr}
\usepackage[parfill]{parskip}
\usepackage[expansion=false]{microtype}
\usepackage{hyperref}
\usepackage{setspace}
\usepackage{needspace}
\usepackage{etoolbox}
\usepackage{graphicx}

\widowpenalty=10000
\clubpenalty=10000
\displaywidowpenalty=10000
\raggedbottom
\sloppy
\hbadness=10000
\hfuzz=2pt

\geometry{top=3.0cm,bottom=3.0cm,left=3.5cm,right=3.5cm}
\setlength{\parskip}{0.65em}
\setlength{\parindent}{0em}
\setlength{\headheight}{14pt}

\definecolor{azulai}{RGB}{30,80,140}
\definecolor{doradoai}{RGB}{140,90,0}
\definecolor{grisai}{RGB}{70,70,70}

\titleformat{\section}{\Large\bfseries\color{azulai}}{}{0em}{}[{\color{azulai}\titlerule[0.5pt]}]
\titlespacing*{\section}{0pt}{1.8em}{0.8em}

\titleformat{\subsection}{\large\bfseries\color{doradoai}}{}{0em}{}
\titlespacing*{\subsection}{0pt}{1.4em}{0.5em}

\titleformat{\subsubsection}{\normalsize\bfseries\color{grisai}}{}{0em}{}
\titlespacing*{\subsubsection}{0pt}{1.0em}{0.3em}

% Evita títulos huérfanos al final de página
\pretocmd{\section}{\Needspace{12\baselineskip}}{}{}
\pretocmd{\subsection}{\Needspace{9\baselineskip}}{}{}
\pretocmd{\subsubsection}{\Needspace{6\baselineskip}}{}{}

\pagestyle{fancy}
\fancyhf{}
\rhead{\textcolor{grisai}{\small Amaya Beatriz — Arquitectura Interna}}
\lhead{\textcolor{grisai}{\small Luna · Casa 4 · Casa 6}}
\cfoot{\textcolor{grisai}{\small\thepage}}
\renewcommand{\headrulewidth}{0.3pt}

\hypersetup{colorlinks=true,linkcolor=azulai,urlcolor=azulai}
\setstretch{1.45}
\tolerance=2500
\emergencystretch=8em

\begin{document}

% ── Portada ──────────────────────────────────────────────────────────────────
\begin{titlepage}
  \centering
  \vspace*{1.5cm}
  {\Huge\bfseries\color{azulai} Luna · Casa 4 · Casa 6}\\[0.5cm]
  {\large\color{grisai} Arquitectura Interna}\\[0.3cm]
  {\small\itshape\color{grisai} Un método para sostener cuerpo, energía y vida con coherencia}\\[2cm]
  {\huge\color{doradoai} Amaya Beatriz}\\[1.5cm]
  {\Large 30/10/1976 \quad 22:10}\\[0.3cm]
  {\Large Zaragoza, España}\\[0.3cm]
  {\normalsize Lat: 41.6916° \quad Lon: -0.9101° \quad UTC+01:00}\\[0.3cm]
  {\normalsize Ascendente: Cáncer 16°08' \quad
    MC: Piscis 25°32'}\\[2cm]

  \begin{tabular}{lll}
    \textbf{Luna:} & Acuario & Casa 8 \\
    \textbf{Casa 4:} & Virgo en cúspide & Regente: Mercurio \\
    \textbf{Casa 6:} & Sagitario en cúspide & Regente: Júpiter \\
  \end{tabular}\\[2cm]

  \vfill
  {\small Generado el 19/05/2026}
\end{titlepage}

\tableofcontents
\clearpage

% ── Datos de referencia ───────────────────────────────────────────────────────
\section{Datos de referencia}

\begin{center}
\begin{tabular}{llll}
  \toprule
  \textbf{Eje} & \textbf{Signo} & \textbf{Posición} & \textbf{Datos} \\
  \midrule
  Luna & Acuario & Casa 8 & 19°05' \\
  Casa 4 (cúspide) & Virgo & --- & Regente: Mercurio \\
  Casa 6 (cúspide) & Sagitario & --- & Regente: Júpiter \\
  Ascendente & Cáncer & --- & 16°08' \\
  Medio Cielo & Piscis & --- & 25°32' \\
  \bottomrule
\end{tabular}
\end{center}

\vspace{0.4cm}
\textbf{Planetas en Casa 4:} Plutón\\
\textbf{Planetas en Casa 6:} Venus, Neptuno

\vspace{0.5cm}
\textbf{Aspectos de la Luna:}

\begin{center}
\begin{tabular}{lllll}
  \toprule
  & \textbf{Asp.} & \textbf{Planeta} & \textbf{Tipo} & \textbf{Orbe} \\
  \midrule
  Luna & ☍ & Saturno & Oposición & 2.9° \\
  Luna & □ & Marte & Cuadratura & 4.0° \\
  Luna & △ & Plutón & Trígono & 6.5° \\
  \bottomrule
\end{tabular}
\end{center}

\vspace{0.8cm}
\begin{center}
\includegraphics[width=0.80\textwidth]{C:/Users/Izas/claude/astrología/pruebas/Amaya_Beatriz_Luna_C4_C6_rueda.png}
\end{center}


\vspace{0.5cm}
\Needspace{5\baselineskip}

% ── Interpretación ────────────────────────────────────────────────────────────
\section{Interpretación — Arquitectura Interna}

\Needspace{8\baselineskip}
\begin{center}
{\small\itshape
No se trata de explicarte. Se trata de mostrar cómo funciona tu regulación interna,\\
dónde puede haber desregulación y qué necesitas para sostenerte.
}
\end{center}

\vspace{0.8cm}
\Needspace{5\baselineskip}

% ── 1. Luna ───────────────────────────────────────────────────────────────────
\subsection{1. Tu Luna — Función emocional}

\subsubsection*{Luna en Acuario — Casa 8}

Tiendes a observar lo que sientes desde cierta distancia.

Muchas veces entiendes emocionalmente lo que te ocurre antes de terminar de habitarlo. Eso puede darte claridad, pero también hacer que tus propias necesidades emocionales lleguen tarde.

Cuando la intensidad emocional supera cierto límite, puede aparecer desconexión repentina. Tu cuerpo suele pedir distancia antes incluso de que seas plenamente consciente de la saturación.

Necesitas espacio interno y margen para procesar lo que sientes a tu manera.

La regulación mejora cuando existe autonomía, tiempo a solas y libertad emocional. En cambio, la invasión emocional, el exceso de intensidad o sentirte sin escapatoria dentro de dinámicas ajenas tienden a generar desconexión y retirada.

Las emociones superficiales rara vez te movilizan demasiado. Lo que realmente te afecta suelen ser las experiencias intensas: la intimidad profunda, las pérdidas, las crisis o los vínculos donde hay algo importante en juego. Tiendes a vivir muchas emociones intensamente en privado y, cuando no existe un espacio seguro para atravesar todo eso, la tensión empieza a acumularse dentro poco a poco. Hay experiencias que puedes sostener en silencio durante muchísimo tiempo sin que casi nadie lo note. Te ayuda sentir confianza profunda, intimidad real y seguridad emocional. En cambio, la traición, las pérdidas emocionales o atravesar demasiada intensidad sin apoyo suele dejar estados internos difíciles de relajar.

Existe tensión entre lo que sientes y la necesidad de mantener control. Cuando emocionalmente algo sube, la respuesta automática suele ser contenerlo o intentar sostenerlo sin mostrar demasiado. El problema es que lo que permanece contenido durante demasiado tiempo termina acumulando presión dentro del cuerpo y del mundo emocional. La sensación de estabilidad aparece cuando la contención deja de ser bloqueo y se convierte en una forma consciente de sostener lo que sientes.

Cuando algo emocionalmente te activa, el cuerpo entra rápidamente en tensión. Si no existe una forma clara de descargar esa energía, puede aparecer irritabilidad, impulsividad o necesidad de confrontación. Muchas veces la reacción responde más al nivel de activación acumulada que a lo que realmente está ocurriendo en ese momento. El aprendizaje no consiste en dejar de sentir intensidad, sino en aprender a darle una salida y un canal más consciente.

Tienes una capacidad muy profunda de transformación emocional. Cuando algo te afecta, puedes atravesarlo con una fuerza interna enorme aunque el proceso no siempre sea sencillo. Existe capacidad real para regenerarte desde dentro y reconstruirte después de experiencias intensas. La emoción aquí no se queda solo en impacto: puede convertirse en comprensión, fuerza y transformación profunda cuando encuentra un cauce adecuado.

\vspace{0.8cm}
\Needspace{5\baselineskip}

% ── 2. Casa 4 ─────────────────────────────────────────────────────────────────
\subsection{2. Tu Casa 4 — Raíz interna}

\subsubsection*{Virgo en la cúspide — Regente: Mercurio}

Cuando aparece sensación de desorden dentro de ti, muchas veces intentas recuperar estabilidad organizando algo fuera. Ordenar, limpiar, resolver tareas o estructurar el entorno puede ayudarte a bajar momentáneamente la tensión interna. La dificultad aparece cuando pasas demasiado tiempo intentando corregir lo externo sin terminar de entrar en lo que realmente te está ocurriendo por dentro. Tu sensación de seguridad aumenta cuando la vida cotidiana tiene cierto orden y funcionalidad. Cuando todo alrededor se vuelve caótico, el cuerpo entra rápidamente en sobrecarga. Te ayuda sentir estructura simple, orden y claridad en lo cotidiano. En cambio, el caos, la sobreexigencia o sentir que nunca llegas a todo suele generar mucha tensión interna.

El regente de la Casa 4 es Mercurio, en Escorpio y en Casa 5. La base interna tiende a organizarse desde la creatividad y la expresión, con la cualidad receptiva y emocional de Escorpio.

Tu base emocional ha atravesado experiencias intensas que dejaron una huella profunda. Hay procesos que transformaron radicalmente tu forma de entender la seguridad emocional y eso te dio muchísima capacidad de resistencia. Al mismo tiempo, ciertas emociones pueden seguir teniendo una intensidad enorme cuando algo las activa. Cuando aparecen temas relacionados con el hogar, la pérdida o el control emocional, la reacción interna puede ser extremadamente poderosa. Tu mundo interior tiene muchísima profundidad y una fuerza emocional difícil de medir desde fuera.

\vspace{0.8cm}
\Needspace{5\baselineskip}

% ── 3. Casa 6 ─────────────────────────────────────────────────────────────────
\subsection{3. Tu Casa 6 — Regulación cotidiana}

\subsubsection*{Sagitario en la cúspide — Regente: Júpiter}

Necesitas sentir que tus hábitos tienen sentido. Las rutinas sostenidas únicamente por obligación suelen agotarte o durar muy poco tiempo. Cuando conectas con algo que realmente te importa, la energía aparece mucho más fácilmente. La dificultad surge cuando desaparece el entusiasmo, porque muchas veces también desaparece toda la estructura cotidiana. Entonces el cuerpo termina pagando el coste de la irregularidad acumulada. Necesitas construir hábitos mínimos capaces de sostenerse incluso en momentos de poca motivación. Te ayudan la dirección, el propósito y sentir movimiento en tu vida. En cambio, la rutina vacía, la desmotivación o la irregularidad constante suele desorganizar rápidamente la energía.

El regente de la Casa 6 es Júpiter, en Tauro y en Casa 11. La regulación cotidiana tiende a organizarse desde la comunidad y los proyectos colectivos, con la cualidad concreta y sostenida de Tauro.

Necesitas cierto bienestar real dentro de la vida cotidiana para sentirte bien físicamente. No suele ayudarte vivir únicamente desde la obligación. La belleza, el disfrute sencillo, la armonía o sentir comodidad en el entorno tienen muchísimo impacto sobre tu energía. Cuando el día a día se vuelve demasiado frío, mecánico o puramente productivo, el cuerpo empieza a vaciarse lentamente. Necesitas que exista algo agradable, humano y habitable dentro de la rutina para poder sostenerla de verdad.

Muchas veces te cuesta notar claramente dónde termina el esfuerzo y dónde empieza el agotamiento. Puedes absorber la carga emocional del entorno sin darte cuenta enseguida y seguir funcionando mientras el cuerpo se va saturando lentamente. Necesitas espacios reales de descanso, silencio y cierre emocional para vaciar todo lo que vas acumulando. Cuando no existe ese vaciado, la fatiga empieza a crecer aunque no haya una causa física evidente. Tu cuerpo necesita mucha más recuperación emocional de la que suele parecer desde fuera.

\vspace{0.8cm}
\Needspace{5\baselineskip}

% ── 4. Integración ────────────────────────────────────────────────────────────
\subsection{4. Integración — Cómo funcionan las tres capas juntas}

Tu Luna en Acuario muestra cómo reaccionas emocionalmente cuando algo te afecta.

Es el primer movimiento: lo que aparece antes de que puedas ordenarlo, explicarlo o decidir qué hacer con ello.

Después, esa emoción busca un lugar donde apoyarse. Ahí entra tu Casa 4, con Virgo en la cúspide: la parte de ti que necesita sentir base, refugio y seguridad interna.

Si esa base está disponible, lo que sientes puede moverse con más facilidad.
Si no lo está, la emoción se queda más tiempo en estado de alerta.

Y lo que no encuentra apoyo interno termina llegando al cuerpo. Ahí entra tu Casa 6, con Sagitario en la cúspide: tus hábitos, tu energía diaria, tu descanso y la forma en que el cuerpo sostiene lo que vives.

Por eso, cuando algo se desordena emocionalmente, no se queda solo en lo emocional. Puede acabar apareciendo como cansancio, pérdida de ritmo, dificultad para descansar o sensación de no poder sostener el día con la misma energía.

En tu caso, las tres capas tienen elementos distintos: Aire en la Luna, Tierra en la Casa 4 y Fuego en la Casa 6.

Esto significa que tu emoción, tu base interna y tu cuerpo no siempre piden lo mismo.

Lo que te ayuda emocionalmente puede no ser lo que tu cuerpo necesita ese día. Lo que te da seguridad interna puede no coincidir con el ritmo que sostiene tu rutina.

Por eso es importante no buscar una única respuesta para todo. A veces una parte de ti necesita expresión, otra necesita refugio y otra necesita descanso, orden o movimiento.

Además, Plutón en la Casa 4 influye directamente en tu sensación de base interna.

Cuando buscas seguridad, refugio o estabilidad emocional profunda, esa energía ya está presente ahí.

Por eso tu forma de sentir hogar, intimidad y raíz no depende solo del signo de la Casa 4. También está marcada por lo que esos planetas activan dentro de ti.

Además, Venus y Neptuno en la Casa 6 influye en cómo tu cuerpo sostiene la vida cotidiana.

No se trata solo de hábitos o rutina. Se trata de cómo organizas tu energía, cómo respondes al cansancio y qué ocurre cuando el cuerpo empieza a pedir ajuste.

Esos planetas muestran qué tipo de fuerza, sensibilidad o tensión entra en tu regulación diaria.

Los aspectos de tensión con Saturno y Marte hacen que el mundo emocional tenga menos margen en ciertos momentos.

Cuando algo te activa, puede haber menos tiempo entre sentir y reaccionar. La emoción puede llegar más intensa, más mezclada o más difícil de ordenar.

Esto no significa que sí o sí vayas a desregularte. Significa que necesitas reconocer antes las señales iniciales, porque cuando la activación ya ha subido mucho, cuesta más volver al centro.

Al mismo tiempo, los aspectos de apoyo con Plutón ofrecen recursos reales.

Cuando consigues encontrar un primer punto de estabilidad, esos aspectos ayudan a sostenerlo.

No eliminan la tensión, pero pueden facilitar que vuelvas a organizarte sin quedarte completamente dentro de lo que se activó.

Lo más difícil no suele ser que una sola parte se altere.

Lo más difícil aparece cuando las tres capas se mueven al mismo tiempo:
la Luna en Acuario, con su forma concreta de sentir;
la Casa 4 en Virgo, con su necesidad de base y seguridad interna;
y la Casa 6 en Sagitario, con el cuerpo intentando sostenerlo todo en la vida diaria.

Cuando esto ocurre, puedes notar que no sabes por dónde empezar: emocionalmente hay activación, internamente falta suelo y el cuerpo empieza a mostrar cansancio, tensión o pérdida de ritmo.

El trabajo no consiste en resolverlo todo a la vez.

Consiste en encontrar la primera capa disponible: una emoción que pueda expresarse, un entorno que pueda darte más seguridad, o una acción corporal sencilla que te ayude a volver poco a poco a la estabilidad.

\vspace{0.8cm}
\Needspace{5\baselineskip}

% ── 5. Sostén y desregulación ─────────────────────────────────────────────────
\subsection{5. Sostén y desregulación}

Tu Luna en Acuario se desregula cuando lo que sientes no encuentra salida.

Puede ser porque no puedes expresarlo, porque no hay espacio para moverte, porque el entorno no lo recibe o porque intentas sostenerlo demasiado tiempo dentro.

No siempre hace falta que ocurra algo muy intenso.

A veces basta con acumular durante demasiado tiempo una emoción que no ha tenido lugar, palabra, movimiento o descanso.

Los aspectos de tensión con Saturno y Marte hacen que ese margen sea más estrecho.

Cuando algo te activa, puede haber menos tiempo entre sentir y reaccionar.

La emoción puede llegar más fuerte, más mezclada o más difícil de ordenar. Por eso es importante reconocer las primeras señales antes de llegar al punto de desborde.

Tu Casa 4 con Virgo pierde estabilidad cuando dejas de sentir base interna.

No siempre tiene que haber un conflicto claro.

A veces basta con que cambie el clima emocional del entorno, con que una relación cercana se tense, con que no puedas descansar de verdad o con que el espacio íntimo deje de sentirse seguro.

Cuando eso ocurre, no solo cambia el ánimo. Cambia la sensación de tener suelo por dentro.

Tu Casa 6 con Sagitario muestra la desregulación en el cuerpo y en la vida cotidiana.

Puede aparecer como cansancio, pérdida de ritmo, tensión física, sueño menos reparador o dificultad para sostener hábitos que normalmente sí puedes sostener.

Muchas veces el cuerpo avisa antes de que puedas explicar qué está pasando.

Cuando el cuerpo empieza a perder continuidad, conviene mirar también qué emoción no ha encontrado salida y qué parte de tu base interna ha dejado de sentirse segura.

Los momentos más difíciles aparecen cuando las tres capas se mueven a la vez.

La Luna en Acuario se activa.
La Casa 4 en Virgo deja de sentirse como base.
Y la Casa 6 en Sagitario empieza a mostrarlo en el cuerpo, el descanso o la rutina.

Ahí puedes sentir que no sabes por dónde empezar.

No se trata de resolverlo todo a la vez.
Se trata de encontrar la primera capa disponible: expresar algo, descansar, ordenar una pequeña rutina, moverte, pedir espacio o recuperar una mínima sensación de seguridad.

Los aspectos de apoyo con Plutón funcionan como recursos cuando consigues recuperar un primer punto de estabilidad.

No eliminan la desregulación, pero pueden ayudarte a reorganizarte antes.

Cuando aparece una mínima base, esos aspectos facilitan que no te quedes dentro de lo que se activó.

\vspace{0.8cm}
\Needspace{5\baselineskip}

% ── 6. Orientación final ──────────────────────────────────────────────────────
\subsection{6. Orientación final}

\subsubsection*{Qué observar}

Empieza observando qué situaciones activan tu Luna en Acuario.

No para corregirte.
No para controlar lo que sientes.

Solo para reconocer antes cómo empieza la activación.

¿Qué ocurre justo antes de sentirte que la situación te sobrepasa?
¿Qué cambia en el entorno, en los vínculos o en el ritmo diario?

Muchas veces la desregulación no empieza en el momento del desborde. Empieza bastante antes: cuando tu Casa 4 en Virgo deja de sentirse segura o cuando tu Casa 6 en Sagitario pierde continuidad y el cuerpo empieza a acumular tensión.

\subsubsection*{Qué puede desregularte}

En tu caso, la desregulación suele seguir una secuencia.

Primero, tu Luna en Acuario busca algo que necesita para estabilizarse y no lo encuentra.

Después, tu Casa 4 deja de funcionar como base interna. Puede ser por tensión emocional, por cambios en el entorno, por exceso de carga o simplemente porque llevas demasiado tiempo sosteniendo más de lo que puedes absorber.

Y finalmente, la Casa 6 en Sagitario empieza a mostrarlo en el cuerpo:
menos energía,
menos descanso,
menos capacidad de sostener hábitos o continuidad.

Cuando las tres capas coinciden, puedes sentir que todo cuesta más de lo habitual aunque desde fuera no parezca estar ocurriendo nada tan grave.

\subsubsection*{Qué te estructura y sostiene}

Lo que te ayuda a regularte no tiene que ser enorme.

De hecho, muchas veces lo más importante son cosas pequeñas que se sostienen en el tiempo.

Una mínima rutina corporal.
Un espacio de descanso real.
Un entorno que se sienta más seguro.
Una forma de expresar lo que ocurre antes de acumularlo demasiado.

Cuando tu Casa 4 siente suelo, tu Luna necesita menos esfuerzo para estabilizarse.

Y cuando tu Casa 6 mantiene cierta continuidad —aunque sea pequeña— el cuerpo puede absorber mejor lo que emocionalmente estás atravesando.

No se trata de no sentir.
Se trata de que lo que sientes no tenga que terminar descargándose siempre en agotamiento, tensión o pérdida de equilibrio.

\subsubsection*{Práctica}

Tu regulación necesita movimiento.

Dedica unos minutos al día a mover el cuerpo conscientemente.

No hace falta intensidad extrema. Hace falta salida.

Caminar rápido, movilizar el cuerpo, sacudir tensión o activar la respiración puede ayudar a que la energía no se quede acumulada internamente.

Cuando el cuerpo puede descargar, la emoción también encuentra más espacio para ordenarse.

\Needspace{8\baselineskip}
\vspace{0.6cm}

\begin{center}
{\small\itshape\color{grisai}
La astrología se usa aquí como lenguaje simbólico de observación, no como definición de la persona.\\
Lo que tienes delante es un mapa de posibilidades, no un destino fijo.
}
\end{center}

\end{document}
